%% Cosinus et sinus d'un réel x (manuel p. 98).
%% M est le point du cercle trigonométrique associé à x (ici x = 0,9 rad) :
%% cos(x) est son abscisse, sin(x) son ordonnée.
\documentclass[tikz,border=4pt]{standalone}
\usepackage{liaisons}
\begin{document}
\begin{tikzpicture}[x=2.2cm,y=2.2cm,
  axe/.style={-{Stealth[length=5pt]}, line width=0.6pt},
  cercle/.style={solideE, line width=1.1pt},
  rappel/.style={line width=0.5pt, dash pattern=on 3pt off 2pt, black!55},
  droit/.style={line width=0.5pt, black!70}]

%% ---------------- coordonnées de M : x = 0,9 rad -----------------------
\def\cosx{0.6216}
\def\sinx{0.7833}

%% ---------------- repère et cercle -------------------------------------
\draw[axe] (-1.28,0) -- (1.32,0);
\draw[axe] (0,-1.28) -- (0,1.32);
\draw[cercle] (0,0) circle (1);
\node[font=\small, text=solideE, anchor=east, inner sep=3pt] at (203:1) {$\mathcal{C}$};
\node[font=\small, anchor=north east, inner sep=2pt] at (-0.03,-0.03) {$O$};
\fill (1,0) circle (1.6pt);
\node[font=\small, anchor=north west, inner sep=3pt] at (1,0) {$I$};
\fill (0,1) circle (1.6pt);
\node[font=\small, anchor=south east, inner sep=3pt] at (0,1) {$J$};

%% ---------------- traits de rappel et angles droits --------------------
\draw[rappel] (\cosx,\sinx) -- (0,\sinx);
\draw[rappel] (\cosx,\sinx) -- (\cosx,0);
\draw[droit] (\cosx-0.1,0) -- (\cosx-0.1,0.1) -- (\cosx,0.1);
\draw[droit] (0.1,\sinx) -- (0.1,\sinx-0.1) -- (0,\sinx-0.1);

%% ---------------- rayon OM, arc x et point M ---------------------------
\draw[solideA, line width=1.2pt] (0,0) -- (\cosx,\sinx);
\draw[solideA, line width=1pt] (0.42,0) arc (0:51.566:0.42);
\node[font=\small, text=solideA, anchor=west, inner sep=2pt] at (24:0.46) {$x$};
\fill[solideA] (\cosx,\sinx) circle (1.8pt);
\node[font=\small, anchor=south west, inner sep=2pt] at (\cosx,\sinx) {$M$};

%% ---------------- projections : cos(x) et sin(x) -----------------------
\fill (\cosx,0) circle (1.6pt);
\node[font=\small, anchor=north, inner sep=4pt] at (\cosx,0) {$\cos(x)$};
\fill (0,\sinx) circle (1.6pt);
\node[font=\small, anchor=east, inner sep=4pt] at (0,\sinx) {$\sin(x)$};
\end{tikzpicture}
\end{document}
